% Extracción quirúrgica del propietario V2 05a: líneas 293--394.
% El resto permanece en este archivo como procedencia, pero no se publica.
\iffalse
\chapter{Macroestructura espectral--polar interna del continuo}
\label{chap:continuo-macro-sp-interna}

\section{Alcance estrictamente interno}

La construcción de este capítulo termina antes de cualquier reconocimiento
externo. Sus únicos datos son el soporte APP de dos hojas, la orientación
TRIT, las extensiones y el carry TPK, el sistema proyectivo de historias, la
medida cilíndrica y el ledger firmado. El resultado es una macroestructura
graduada con transporte, expresión de cambio de hoja, lector firmado y
telescopía no autónoma. No se utiliza una función exterior para definir
ninguno de sus mapas.

\section{Dominio superviviente de la rama}

Para \(x_k\in\mathcal A_k\), sea
\begin{equation}
\label{eq:continuo-sp-msect-n}
 \operatorname{Msect}^{\mathrm{sp},(N)}_k(x_k)
 =\left\{(x_j)_{j=k}^{N}:
 \tau_{j+1,j}x_{j+1}=x_j,
 \ (x_j,x_{j+9})\in\Gamma_{9,j},
 \ \text{y se conserva el haz }\{\Sigma,\Pi\}\right\}.
\end{equation}
Definimos
\begin{equation}
\label{eq:continuo-sp-dominio-superviviente}
 \mathcal C_k^{\mathrm{sp}}
 =\left\{x_k:
 \operatorname{Msect}^{\mathrm{sp},(N)}_k(x_k)\neq\varnothing
 \text{ para todo }N\geq k\right\},
 \qquad
 \mathcal C_{\rm cont}^{\mathrm{sp}}
 =\varprojlim_k\mathcal C_k^{\mathrm{sp}}.
\end{equation}
El umbral TRIT no se excluye de este dominio. El haz conserva las dos hojas
aritméticas aunque una carta visible no publique ninguna de ellas en un
evento concreto.

\section{Paquete de trece coordenadas}

Para \(x_k\in\mathcal C_k^{\mathrm{sp}}\), el paquete dependiente es
\begin{equation}
\label{eq:continuo-sp-d13}
\begin{aligned}
 D_{\mathrm{sp},k}(x_k)=\bigl(&
 \mathcal F_{\mathrm{sp},k}(x_k),
 \operatorname{Ext}_{\Gamma,\mathrm{sp},k}(x_k),
 \operatorname{Rel}_{\mathrm{sp},k}(x_k),
 \mathcal N_{\mathrm{sp},k}(x_k),
 \operatorname{or}_{\mathrm{sp},k}(x_k),\\
 &\operatorname{Carr}_{\mathrm{sp},k}(x_k),
 \operatorname{Mem}_{\mathrm{sp},k}(x_k),
 \partial_{\mathrm{sp},k}(x_k),
 \gamma_{\mathrm{sp},k}(x_k),
 \operatorname{Msect}_{\mathrm{sp},k}(x_k),\\
 &\operatorname{Surv}_{\mathrm{sp},k}(x_k),
 \operatorname{Term}_{\mathrm{sp},k}(x_k),
 \mathcal L^{\rm sgn}_{\mathrm{sp},k}(x_k)\bigr).
\end{aligned}
\end{equation}
Sus componentes se determinan correlativamente como sigue.

La fibra \(\mathcal F_{\mathrm{sp},k}\) contiene el estado TPK y el haz
ordenado \(\{\Sigma,\Pi\}\). La extensión es la restricción completa de
\(\operatorname{Ext}_\Gamma\) a historias que satisfacen
\eqref{eq:continuo-sp-dominio-superviviente}. Las relaciones contienen las
leyes de identidad y composición, el lift residuo--cociente, el cociclo de
carry, la compatibilidad de frontera y la naturalidad del truncamiento. La
coordenada numérica contiene profundidad, fase, residuos, cocientes y firma
interna. La orientación almacena por separado hoja y orientación TRIT. La
memoria contiene el cociente, el ledger y el incremento de vuelta. La
frontera y la ruta permanecen sin contraer. La multisección es el sistema
completo de prolongaciones de \eqref{eq:continuo-sp-msect-n}. La
supervivencia es la pertenencia estable de
\eqref{eq:continuo-sp-dominio-superviviente}. El terminal y el lector se
construyen en las secciones siguientes.

Los truncamientos de todas las coordenadas forman un mapa
\(u^{\mathrm{sp}}_{\ell,k}\) que satisface
\begin{equation}
\label{eq:continuo-sp-d-naturalidad}
 u^{\mathrm{sp}}_{\ell,k}D_{\mathrm{sp},\ell}
 =D_{\mathrm{sp},k}\tau_{\ell,k}.
\end{equation}

\section{Gráfica de la restricción y sección dependiente}

Se define
\begin{equation}
\label{eq:continuo-sp-g-graph}
 \mathcal G_{\mathrm{sp},k}
 =\operatorname{Graph}(D_{\mathrm{sp},k})
 =\{(x_k,D_{\mathrm{sp},k}x_k):
     x_k\in\mathcal C_k^{\mathrm{sp}}\}.
\end{equation}
La restricción total es
\begin{equation}
\label{eq:continuo-sp-res-graph}
 \operatorname{Res}_{\mathrm{sp}}=s_{\mathrm{sp}}:
 \mathcal C_{\rm cont}^{\mathrm{sp}}
 \longrightarrow\mathcal G_{\mathrm{sp}},
 \qquad
 s_{\mathrm{sp}}(x)=(x,D_{\mathrm{sp}}x).
\end{equation}
La proyección \(\pi_1(x,d)=x\) verifica
\begin{equation}
\label{eq:continuo-sp-res-inversa}
 \pi_1s_{\mathrm{sp}}=I,
 \qquad
 s_{\mathrm{sp}}\pi_1=I_{\mathcal G_{\mathrm{sp}}}.
\end{equation}

Sea \(r_{\mathrm{sp},k}(x_k)\) la multisección de hojas y rutas que ya
pertenece a \(D_{\mathrm{sp},k}(x_k)\). Se define
\begin{equation}
\label{eq:continuo-sp-x-dependiente}
 X_{\chi,\mathrm{sp},k}
 =\{(r_{\mathrm{sp},k}x_k,x_k,D_{\mathrm{sp},k}x_k):
     x_k\in\mathcal C_k^{\mathrm{sp}}\},
\end{equation}
\begin{equation}
\label{eq:continuo-sp-pr-x}
 \operatorname{pr}_{X,\mathrm{sp}}(x,D_{\mathrm{sp}}x)
 =(r_{\mathrm{sp}}x,x,D_{\mathrm{sp}}x).
\end{equation}
Olvidar la primera coordenada de
\eqref{eq:continuo-sp-x-dependiente} es la inversa de
\(\operatorname{pr}_{X,\mathrm{sp}}\).

\section{Haz de hojas y medida proyectiva}

Sea
\begin{equation}
\label{eq:continuo-sp-haz-hojas}
 \widetilde{\mathcal C}_k^{\mathrm{sp}}
 =\{(x_k,h):x_k\in\mathcal C_k^{\mathrm{sp}},
              h\in\{\Sigma,\Pi\}\}.
\end{equation}
La medida proyectiva del continuo induce probabilidades
\(\widetilde\mu_k\) sobre este haz. Se restringe al soporte positivo; para
\(z\) en ese soporte,
\begin{equation}
\label{eq:continuo-sp-consistencia-medida}
 \widetilde\mu_k(z)
 =\sum_{\widetilde\tau_{k+1,k}(z')=z}
  \widetilde\mu_{k+1}(z').
\end{equation}
La linealización de nivel es
\begin{equation}
\label{eq:continuo-sp-hilbert-nivel}
 \mathcal H_k^{\mathrm{sp}}
 =\ell^2(\operatorname{supp}\widetilde\mu_k),
\end{equation}
con base ortonormal \(e_z\).

\section{Transporte de todas las extensiones y carry}

Para \(z'\) que prolonga \(z\), definimos el peso condicional positivo
\begin{equation}
\label{eq:continuo-sp-peso-condicional}
 w_k(z'\mid z)
 =\left(\frac{\widetilde\mu_{k+1}(z')}
               {\widetilde\mu_k(z)}\right)^{1/2}.
\end{equation}
El transporte es
\begin{equation}
\label{eq:continuo-sp-u-accion}
 U_ke_z
 =\sum_{\substack{z':\widetilde\tau_{k+1,k}(z')=z\\
            (z,z')\in\operatorname{Ext}_{\Gamma,\mathrm{sp},k}}}
 w_k(z'\mid z)e_{z'}.
\end{equation}
La suma contiene todas las extensiones admisibles. No depende de un selector.

Por \eqref{eq:continuo-sp-consistencia-medida},
\[
 \sum_{z':\widetilde\tau(z')=z}|w_k(z'\mid z)|^2=1.
\]
Los conjuntos de descendientes de dos padres distintos son disjuntos. En
consecuencia,
\begin{equation}
\label{eq:continuo-sp-u-isometria}
 U_k^*U_k=I_{\mathcal H_k^{\mathrm{sp}}}.
\end{equation}

Cada base \(e_{z'}\) incluye el carry de la extensión correspondiente. Para
rutas componibles se conservan
\begin{align}
 \kappa(\gamma_2\gamma_1)
 &=\kappa(\gamma_2)+
   \mathsf C(\gamma_2)\kappa(\gamma_1),
 \label{eq:continuo-sp-carry-nativo}\\
 \operatorname{Carr}(\gamma_2\gamma_1)
 &=\operatorname{Carr}(\gamma_2)H(\gamma_1)
  +Y(\gamma_2)\operatorname{Carr}(\gamma_1).
 \label{eq:continuo-sp-carry-operador}
\end{align}
Los pesos condicionales se multiplican a lo largo de la única cadena de
truncamientos de un descendiente. Las ecuaciones
\eqref{eq:continuo-sp-carry-nativo}--
\eqref{eq:continuo-sp-carry-operador} garantizan que las etiquetas de carry
coincidan. Por tanto,
\begin{equation}
\label{eq:continuo-sp-u-composicion}
 U_{\ell,k}=U_{\ell-1}\cdots U_k.
\end{equation}

\section{Graduación, proyectores y descomposición del transporte}

La graduación de hoja es la involución
\begin{equation}
\label{eq:continuo-sp-sigma}
 \sigma_ke_{(x,\Sigma)}=+e_{(x,\Sigma)},
 \qquad
 \sigma_ke_{(x,\Pi)}=-e_{(x,\Pi)}.
\end{equation}
Así,
\[
 \sigma_k^*=\sigma_k,
 \qquad
 \sigma_k^2=I.
\]
Los proyectores internos son
\begin{equation}
\label{eq:continuo-sp-p-q}
 P_k=\frac{I+\sigma_k}{2},
 \qquad
 Q_k=\frac{I-\sigma_k}{2}.
\end{equation}
Satisfacen
\begin{equation}
\label{eq:continuo-sp-p-q-relaciones}
 P_k^2=P_k,
 \quad Q_k^2=Q_k,
 \quad P_kQ_k=0,
 \quad P_k+Q_k=I.
\end{equation}

Separamos la parte del transporte que conserva hoja de la que la cambia:
\begin{align}
 A_k&=P_{k+1}U_kP_k+Q_{k+1}U_kQ_k,
 \label{eq:continuo-sp-a}\\
 \eta_k&=Q_{k+1}U_kP_k+P_{k+1}U_kQ_k.
 \label{eq:continuo-sp-eta}
\end{align}
Entonces
\begin{equation}
\label{eq:continuo-sp-u-a-eta}
 U_k=A_k+\eta_k,
\end{equation}
y la graduación verifica
\begin{equation}
\label{eq:continuo-sp-paridad-a-eta}
 \sigma_{k+1}A_k=A_k\sigma_k,
 \qquad
 \sigma_{k+1}\eta_k=-\eta_k\sigma_k.
\end{equation}

\begin{proposition}[Balance local de hojas]
\label{prop:continuo-sp-balance-local}
Para cada nivel,
\begin{equation}
\label{eq:continuo-sp-balance-local}
 A_k^*A_k+\eta_k^*\eta_k=I.
\end{equation}
\end{proposition}

\begin{proof}
La ortogonalidad de \(P_{k+1}\) y \(Q_{k+1}\) da
\begin{align*}
 A_k^*A_k
 &=P_kU_k^*P_{k+1}U_kP_k
   +Q_kU_k^*Q_{k+1}U_kQ_k,\\
 \eta_k^*\eta_k
 &=P_kU_k^*Q_{k+1}U_kP_k
   +Q_kU_k^*P_{k+1}U_kQ_k.
\end{align*}
Sumando y usando \(P_{k+1}+Q_{k+1}=I\),
\[
 A_k^*A_k+\eta_k^*\eta_k
 =P_kU_k^*U_kP_k+Q_kU_k^*U_kQ_k
 =P_k+Q_k=I.
\]
\end{proof}

El operador \(\eta_k\) registra el abandono del canal que conserva hoja. No
es el carry entero: las coordenadas
\eqref{eq:continuo-sp-carry-nativo}--
\eqref{eq:continuo-sp-carry-operador} permanecen en el paquete
\eqref{eq:continuo-sp-d13}.

\fi
\section{Lector firmado con registro genealógico}

Para \(v\in\mathcal H_k^{\mathrm{sp}}\), la polarización firmada es
\begin{equation}
\label{eq:continuo-sp-polarizacion-firmada}
 \operatorname{pol}_k(v)
 =\langle v,\sigma_kv\rangle
 =\|P_kv\|^2-\|Q_kv\|^2.
\end{equation}
El lector especializado conserva además la fuente:
\begin{equation}
\label{eq:continuo-sp-lector-firmado}
 \mathcal L^{\rm sgn}_{\mathrm{sp},k}(x_k;v)
 =\bigl(
   \Lambda(\gamma_{x_k}),
   \Sigma_{\rm reg}(\Lambda(\gamma_{x_k})),
   \|P_kv\|^2,
   \|Q_kv\|^2,
   \operatorname{pol}_k(v)\bigr).
\end{equation}
La orientación TRIT se conserva en \(\Lambda(\gamma_{x_k})\); no se
identifica con el signo de hoja. A horizonte \(N\), el registro guarda la
familia de lectores para todos los niveles \(k\leq N\). Truncar significa
eliminar únicamente las entradas posteriores.

\section{Producto no autónomo y terminal}

Se define
\begin{equation}
\label{eq:continuo-sp-f-producto}
 F_{0,0}=I_{\mathcal H_0^{\mathrm{sp}}},
 \qquad
 F_{0,k}=A_{k-1}\cdots A_1A_0
 \quad(k\geq1).
\end{equation}
Por tanto,
\begin{equation}
\label{eq:continuo-sp-f-recursion}
 F_{0,k+1}=A_kF_{0,k}.
\end{equation}
El término publicado al abandonar el canal en el paso \(k\) y el terminal a
horizonte \(N\) son
\begin{equation}
\label{eq:continuo-sp-defecto-terminal}
 \mathcal D_k=F_{0,k}^*\eta_k^*\eta_kF_{0,k},
 \qquad
 \mathcal T_N=F_{0,N}^*F_{0,N}.
\end{equation}

\begin{theorem}[Telescopía espectral--polar no autónoma]
\label{thm:continuo-sp-telescopia-no-autonoma}
Para todo horizonte finito \(N\),
\begin{equation}
\label{eq:continuo-sp-telescopia-no-autonoma}
 \boxed{
 I=\sum_{k=0}^{N-1}
 F_{0,k}^*\eta_k^*\eta_kF_{0,k}
 +F_{0,N}^*F_{0,N}.}
\end{equation}
Además,
\[
 0\preceq\mathcal T_{N+1}
 \preceq\mathcal T_N\preceq I,
\]
de modo que existe el límite fuerte positivo
\[
 \mathcal T_\infty
 =\operatorname{s-lim}_{N\to\infty}\mathcal T_N
\]
y
\begin{equation}
\label{eq:continuo-sp-telescopia-infinita}
 I=\sum_{k=0}^{\infty}
 F_{0,k}^*\eta_k^*\eta_kF_{0,k}
 +\mathcal T_\infty
\end{equation}
en topología fuerte. El teorema no afirma que
\(\mathcal T_\infty\) sea nulo.
\end{theorem}

\begin{proof}
Multiplicando \eqref{eq:pdfv2-cor-balance-hojas} por
\(F_{0,k}^*\) y \(F_{0,k}\), y usando
\eqref{eq:continuo-sp-f-recursion}, se obtiene
\[
 F_{0,k}^*F_{0,k}
 =F_{0,k+1}^*F_{0,k+1}
  +F_{0,k}^*\eta_k^*\eta_kF_{0,k}.
\]
La suma desde \(k=0\) hasta \(N-1\) cancela todos los términos
intermedios y deja precisamente
\eqref{eq:continuo-sp-telescopia-no-autonoma}. Como
\(A_k^*A_k\preceq I\),
\[
 \mathcal T_{N+1}
 =F_{0,N}^*A_N^*A_NF_{0,N}
 \preceq F_{0,N}^*F_{0,N}=\mathcal T_N.
\]
La monotonía y la cota positiva proporcionan el límite fuerte. Pasar al
límite en la identidad finita da
\eqref{eq:continuo-sp-telescopia-infinita}.
\end{proof}

\iffalse
\section{Ensamblador como gráfica}

Para
\[
 z_N=(r_{\mathrm{sp},N}x_N,x_N,D_{\mathrm{sp},N}x_N)
 \in X_{\chi,\mathrm{sp},N},
\]
se define el ensamblador material
\begin{equation}
\label{eq:continuo-sp-a-accion}
\begin{aligned}
 a_{\mathrm{sp},N}(z_N)=\bigl(&
 (\mathcal H_k^{\mathrm{sp}},\sigma_k,P_k,Q_k)_{k\leq N},
 (U_k,A_k,\eta_k)_{k<N},\\
 &(F_{0,k},\mathcal D_k,\mathcal T_k,
   \mathcal L^{\rm sgn}_{\mathrm{sp},k})_{k\leq N}\bigr).
\end{aligned}
\end{equation}
La macroestructura es
\begin{equation}
\label{eq:continuo-sp-m-graph}
 \mathfrak M_{\mathrm{sp},N}
 =\operatorname{Graph}(a_{\mathrm{sp},N}),
 \qquad
 \mathcal A_{\mathrm{sp},N}(z_N)
 =(z_N,a_{\mathrm{sp},N}z_N).
\end{equation}
La proyección a \(z_N\) es la inversa de
\(\mathcal A_{\mathrm{sp},N}\). El truncamiento elimina las entradas cuyo
índice supera el horizonte nuevo y verifica
\begin{equation}
\label{eq:continuo-sp-a-naturalidad}
 m^{\mathrm{sp}}_{\ell,k}a_{\mathrm{sp},\ell}
 =a_{\mathrm{sp},k}\xi^{\mathrm{sp}}_{\ell,k}.
\end{equation}

\section{Publicador como gráfica}

El publicador no contrae la macroestructura a un escalar. Devuelve el
certificado completo
\begin{equation}
\label{eq:continuo-sp-p-accion}
 p_{\mathrm{sp},N}(m_N)
 =\left(
 \mathcal L^{\rm sgn}_{\mathrm{sp},k},
 \mathcal D_k,
 \mathcal T_k,
 I=\sum_{j=0}^{k-1}\mathcal D_j+\mathcal T_k
 \right)_{0\leq k\leq N}.
\end{equation}
Se define
\begin{equation}
\label{eq:continuo-sp-c-graph}
 \mathcal C_{\mathrm{sp},N}
 =\operatorname{Graph}(p_{\mathrm{sp},N}),
 \qquad
 \mathcal P_{\mathrm{sp},N}(m_N)
 =(m_N,p_{\mathrm{sp},N}m_N).
\end{equation}
La proyección a \(m_N\) es su inversa. Truncar el certificado elimina sólo
las identidades posteriores, por lo que
\begin{equation}
\label{eq:continuo-sp-p-naturalidad}
 c^{\mathrm{sp}}_{\ell,k}p_{\mathrm{sp},\ell}
 =p_{\mathrm{sp},k}m^{\mathrm{sp}}_{\ell,k}.
\end{equation}

\section{Teorema de la cadena interna completa}

\begin{theorem}[Macroestructura espectral--polar generada por el continuo]
\label{thm:continuo-sp-cadena-interna}
Las acciones
\eqref{eq:continuo-sp-res-graph},
\eqref{eq:continuo-sp-pr-x},
\eqref{eq:continuo-sp-m-graph} y
\eqref{eq:continuo-sp-c-graph} forman la cadena natural
\begin{equation}
\label{eq:continuo-sp-cadena-interna}
 \mathcal C_{\rm cont}^{\rm disc}
 \supset\mathcal C_{\rm cont}^{\mathrm{sp}}
 \xrightarrow[\cong]{\operatorname{Res}_{\mathrm{sp}}}
 \mathcal G_{\mathrm{sp}}
 \xrightarrow[\cong]{\operatorname{pr}_{X,\mathrm{sp}}}
 X_{\chi,\mathrm{sp}}
 \xrightarrow[\cong]{\mathcal A_{\mathrm{sp}}}
 \mathfrak M_{\mathrm{sp}}
 \xrightarrow[\cong]{\mathcal P_{\mathrm{sp}}}
 \mathcal C_{\mathrm{sp}}.
\end{equation}
Su contenido publicado es el haz de dos hojas, la graduación
\(\sigma_k\), los proyectores \(P_k,Q_k\), el transporte total \(U_k\),
la separación \(U_k=A_k+\eta_k\), el carry y la memoria, el lector firmado
y la telescopía
\eqref{eq:continuo-sp-telescopia-no-autonoma} con terminal explícito. Cada
objeto conserva una proyección canónica hacia todos sus antecedentes.
\end{theorem}

\begin{proof}
Las identidades
\eqref{eq:continuo-sp-res-inversa} prueban la primera equivalencia. La
definición dependiente \eqref{eq:continuo-sp-x-dependiente} prueba la
segunda. Las gráficas
\eqref{eq:continuo-sp-m-graph} y
\eqref{eq:continuo-sp-c-graph} prueban las dos restantes. La naturalidad de
los cuatro pasos es
\eqref{eq:continuo-sp-d-naturalidad},
\eqref{eq:continuo-sp-a-naturalidad} y
\eqref{eq:continuo-sp-p-naturalidad}, junto con la naturalidad de
\(r_{\mathrm{sp},k}\). El contenido operatorial queda probado por
\eqref{eq:continuo-sp-u-isometria}, el
\cref{prop:continuo-sp-balance-local} y el
\cref{thm:continuo-sp-telescopia-no-autonoma}. Todas las extensiones aparecen
en \eqref{eq:continuo-sp-u-accion}; por ello no interviene una elección de
ruta.
\end{proof}

\section{Retorno de fase y evolución del estado}

La rama conserva expresamente
\begin{equation}
\label{eq:continuo-sp-no-reset}
 g(k+9)=g(k),
 \qquad
 x_{k+9}\neq x_k\quad\text{en general}.
\end{equation}
El segundo miembro registra el incremento de memoria y los posibles cambios
de hoja, carry, cilindro y frontera. Ni \(\sigma_k\), ni los proyectores, ni
el lector firmado contraen esta diferencia.

\section{Procedencia}

Las dos hojas APP, el estado TRIT enriquecido, las fibras y extensiones TPK,
el carry, el retorno sin reinicialización, la supervivencia y el ledger son
\texttt{ARQUITECTURA\_AUTORAL\_PREEXISTENTE}. La obligación de conservar el
terminal finito es \texttt{RESULTADO\_RECUPERADO}. El haz explícito de hojas,
la involución \(\sigma_k\), los proyectores, la linealización que suma todas
las extensiones, la descomposición \(U_k=A_k+\eta_k\), el producto no
autónomo \(F_{0,k}\) y las gráficas dependientes son
\texttt{FORMALIZACION\_NUEVA}. El
\cref{thm:continuo-sp-telescopia-no-autonoma,thm:continuo-sp-cadena-interna}
constituye el \texttt{CERTIFICADO\_NUEVO} de esa formalización.

\section{Falsadores internos}

\begin{proposition}[Falsadores de la macroestructura]
\label{prop:continuo-sp-falsadores}
La conclusión del
\cref{thm:continuo-sp-cadena-interna} no es transferible si ocurre cualquiera
de las siguientes sustituciones:
\begin{enumerate}
 \item elegir una prolongación en
       \eqref{eq:continuo-sp-u-accion} y borrar las restantes;
 \item sustituir \(\mathcal G_{\mathrm{sp}}\) por la imagen exterior de
       \(D_{\mathrm{sp}}\);
 \item borrar \((x,D_{\mathrm{sp}}x)\) al formar
       \(X_{\chi,\mathrm{sp}}\);
 \item identificar la hoja con la orientación TRIT o borrar el umbral;
 \item romper la consistencia
       \eqref{eq:continuo-sp-consistencia-medida};
 \item borrar el carry o incumplir
       \eqref{eq:continuo-sp-carry-nativo}--
       \eqref{eq:continuo-sp-carry-operador};
 \item identificar \(\eta_k\) con todo el carry;
 \item reemplazar la familia \((A_k)_k\) por un solo operador sin demostrar
       estacionariedad;
 \item suprimir \(F_{0,N}^*F_{0,N}\) en un horizonte finito;
 \item afirmar \(\mathcal T_\infty=0\) sin una prueba adicional;
 \item reconstruir una historia desde la polarización firmada;
 \item usar \(g(k+9)=g(k)\) para afirmar \(x_{k+9}=x_k\);
 \item introducir un dominio, una función o un criterio exterior para
       definir cualquiera de las flechas internas;
 \item introducir una salida posterior como antecedente de la rama.
\end{enumerate}
\end{proposition}
\fi
